\chapter{Insurance and Risk Sharing}\label{ch:b2:insurance-and-risk-sharing}

A household cannot change the weather, a factory fire or a lay-off, but it can change who bears the
loss. Insurance contracts, mutual societies, family transfers, savings and public programmes all move
income from the states of the world where it is plentiful to those where it is scarce. If those
arrangements worked perfectly, a country's consumption would not rise and fall with its own output:
its households would share their risks with the rest of the world, and only the world's common
fluctuations would remain. The chart tests that idea on 123 countries.
% figdata: summary (n 123)

\begin{openingchart}{Year-to-year volatility of household consumption per person against that of GDP
  per person, 1991--2023: one point per country with at least 25 years of both series. Standard
  deviations of the annual growth rates, in percentage points; times-ten scales on both axes. Points
  above the dotted line are countries whose consumption moved more than their output.}{World Bank,
  World Development Indicators: growth of household final consumption expenditure per person
  (NE.CON.PRVT.PC.KD.ZG) and of GDP per person (NY.GDP.PCAP.KD.ZG), constant prices}
\begin{tikzpicture}
\begin{loglogaxis}[omscatter, xmin=1, xmax=25, ymin=1, ymax=25, log basis x=10, log basis y=10,
    log ticks with fixed point, xtick={1,10}, ytick={1,10}, minor tick num=0,
    xlabel={Volatility of GDP growth (points)}, ylabel={Volatility of consumption growth (points)},
    height=0.62\linewidth, width=0.72\linewidth]
  \addplot+[mark size=1.3pt] table[x=sd_gdp, y=sd_cons]
    {figdata/bachelor-2/02-insurance-and-risk-sharing/out/volatility.csv};
  \draw[black!70, densely dotted, thick] (axis cs:1,1) -- (axis cs:25,25);
  \node[font=\footnotesize, anchor=north west] at (axis cs:12,11) {equal volatility};
\end{loglogaxis}
\end{tikzpicture}
\end{openingchart}

Full sharing of risk predicts points well below the dotted line: consumption should move with the
world, and less than output. Instead, 73 per cent of the countries lie above the line, and the median
% figdata: summary (share_ratio_above_1 73.2, median_ratio 1.18)
country's consumption growth is 1.18 times as volatile as its GDP growth: 1.53 times among low-income
% figdata: summary (median_ratio_low_income 1.53, median_ratio_high_income 1.12)
countries, 1.12 times among high-income ones. This chapter builds the theory of insurance and risk sharing that makes the
prediction. It then tests that theory twice: on villages, where it does far better, and on countries,
where it does far worse.

\begin{recall}
From the previous chapter: expected utility, \omterm{def:b2:expected-utility:risk-aversion}{risk aversion} as the concavity of the utility of money,
the \omterm{def:b2:expected-utility:certainty-equivalent}{certainty equivalent} and the \omterm{def:b2:expected-utility:risk-premium}{risk premium}, CARA and CRRA utilities, and the small-risk
approximation of the premium, $\tfrac12\mathcal{A}\sigma^2$. From the school volumes: pooling many
independent risks makes the average loss predictable. From mathematics: the law of large numbers.
\end{recall}

\begin{notation}[Symbols of this chapter]\label{not:b2:insurance-and-risk-sharing:symbols}
As in \cref{ch:b2:expected-utility}: money $z$, claim probability $\mu$, cover $\chi$ bought at $\psi$
per unit. A household with wealth $z_0$ faces a loss $\Lambda$; $\eta$ is a \omterm{def:b2:insurance-and-risk-sharing:loading}{loading} and $d$ a
\omterm{def:b2:insurance-and-risk-sharing:deductible}{deductible}. States of the world are $s$, with probabilities $\mu_s$; household $i$'s endowment is
$\omega_i(s)$ and the aggregate endowment $\Omega(s)$. In the tests, household $i$'s consumption and
income in year $t$ are $c_{it}$ and $m_{it}$. A portfolio puts the share $\theta$ of wealth in an
asset whose return exceeds the safe real rate $r$ by $\tilde\xi$, with mean $\bar\xi$ and variance
$\sigma^2$.
\end{notation}

\section{The demand for insurance}\label{sec:b2:insurance-and-risk-sharing:demand}

A household with wealth $z_0$ loses $\Lambda$ with probability $\mu$. An insurer sells cover: for each
unit of cover $\chi$, the household pays $\psi$ in every state and receives $1$ if the loss occurs.
Its final wealth is
\[
z_G = z_0 - \psi\chi \quad\text{(no loss)}, \qquad z_B = z_0 - \Lambda + (1 - \psi)\chi \quad
\text{(loss)}.
\]

\begin{definition}[Actuarially fair premium]\label{def:b2:insurance-and-risk-sharing:fair-premium}
A premium is an \emph{actuarially fair premium}\index{actuarially fair premium} if it equals the
expected payment of the contract: $\psi\chi = \mu\chi$, that is, a price $\psi = \mu$ per unit of
cover. The insurer then breaks even on average.
\end{definition}

\begin{definition}[Loading]\label{def:b2:insurance-and-risk-sharing:loading}
The \emph{loading}\index{loading} of a premium is its excess over the fair premium, as a share of the
fair premium: $\eta = \psi/\mu - 1$. It pays for the insurer's costs of selling contracts and settling
claims, and for its capital.
\end{definition}

\begin{definition}[Deductible]\label{def:b2:insurance-and-risk-sharing:deductible}
A contract has a \emph{deductible}\index{deductible} $d$ if it pays $\max(x - d, 0)$ on a loss of
size $x$: the household bears every loss up to $d$, and the insurer the rest.
\end{definition}

\begin{proposition}[Full cover at a fair price]\label{prop:b2:insurance-and-risk-sharing:mossin}
Let $u$ be strictly concave and differentiable, and let the household choose $\chi \in [0, \Lambda]$
to maximise $(1 - \mu)u(z_G) + \mu u(z_B)$.
\begin{enumerate}
  \item At a fair price ($\psi = \mu$), it buys full cover, $\chi = \Lambda$, and its wealth is $z_0 -
        \mu\Lambda$ in both states.
  \item With a positive \omterm{def:b2:insurance-and-risk-sharing:loading}{loading} ($\psi > \mu$), it buys less than full cover, and nothing at all when
        $(1 - \mu)\psi\,u'(z_0) \ge \mu(1 - \psi)\,u'(z_0 - \Lambda)$.
\end{enumerate}
\end{proposition}

\begin{proof}
Expected utility $V(\chi)$ is strictly concave in $\chi$, and
\[
V'(\chi) = -(1 - \mu)\psi\,u'(z_G) + \mu(1 - \psi)\,u'(z_B).
\]
(1) With $\psi = \mu$, $V'(\chi) = \mu(1 - \mu)\bigl(u'(z_B) - u'(z_G)\bigr)$. For $\chi < \Lambda$,
$z_B < z_G$, so $u'(z_B) > u'(z_G)$ and $V' > 0$: more cover is better up to $\chi = \Lambda$, where
$z_B = z_G$. (2) At $\chi = \Lambda$, $z_B = z_G$ and $V'(\Lambda) = u'(z_G)(\mu - \psi) < 0$, so the
optimum is below $\Lambda$. It is $0$ when $V'(0) \le 0$, which is the stated inequality.
\end{proof}

\begin{example}[A loading of a quarter]\label{ex:b2:insurance-and-risk-sharing:loading}
Let $u = \ln$, $z_0 = \num{100000}$, $\Lambda = \num{40000}$, $\mu = 0.05$. At the fair price the
household buys \num{40000} of cover for a premium of \num{2000}. With a \omterm{def:b2:insurance-and-risk-sharing:loading}{loading} of 25 per cent,
$\psi = 0.0625$, and $V'(\chi) = 0$ gives
\[
\chi^{*} = \frac{\mu(1 - \psi)z_0 - (1 - \mu)\psi(z_0 - \Lambda)}{\psi(1 - \psi)} =
\frac{\num{4687.5} - \num{3562.5}}{\num{0.05859375}} = \num{19200},
\]
under half of the loss. A \omterm{def:b2:insurance-and-risk-sharing:loading}{loading} of a quarter halves the cover.
\end{example}

The \omterm{def:b2:insurance-and-risk-sharing:deductible}{deductible} is the form most contracts take, and expected utility explains why.

\begin{proposition}[Deductibles (Arrow)]\label{prop:b2:insurance-and-risk-sharing:arrow}
Let losses take many values, and let the premium be the expected payment times a fixed factor $1 +
\eta$. Among all contracts with the same premium whose payment is between $0$ and the loss, a
risk-averse household prefers a contract with a straight \omterm{def:b2:insurance-and-risk-sharing:deductible}{deductible}.
\end{proposition}
\admitted

The intuition: a \omterm{def:b2:insurance-and-risk-sharing:deductible}{deductible} concentrates the cover on the large losses, where marginal utility is
highest, and leaves the small ones, which cost little utility, with the household.

\begin{model}{Demand for insurance}
\assumes A household maximises expected utility of wealth, knows the probability of its loss, and
buys cover at a price set before the loss.
\predicts Full cover at a fair price; partial cover, falling as the \omterm{def:b2:insurance-and-risk-sharing:loading}{loading} rises, when the price
exceeds the expected payment; \omterm{def:b2:insurance-and-risk-sharing:deductible}{deductibles} rather than other partial contracts.
\tested \omterm{def:b2:insurance-and-risk-sharing:deductible}{Deductible} and cover choices in home and flood insurance (Collier et al.\ 2018: most households
take low \omterm{def:b2:insurance-and-risk-sharing:deductible}{deductibles} and full cover at loads well above one); purchases of rainfall insurance at
randomised prices (Karlan et al.\ 2014: fewer buyers at higher prices).
\breaks When the buyer knows its risk better than the insurer (\cref{ch:b2:adverse-selection}), when
cover changes the care the buyer takes (\cref{ch:b2:moral-hazard}), and in the size of the demand for
small-stakes cover (\cref{ch:b2:expected-utility}).
\end{model}

\section{State-contingent goods}\label{sec:b2:insurance-and-risk-sharing:contingent}

Insurance is a trade between states of the world. That idea turns the theory of choice under risk into
ordinary consumer theory.

\begin{definition}[State of the world]\label{def:b2:insurance-and-risk-sharing:state}
A \emph{state of the world}\index{state of the world} is a complete description of whatever is
uncertain and relevant to the choice: whether the house burns, how much rain falls, which firms
succeed. The states $s = 1, \dots, S$ are mutually exclusive, and exactly one occurs.
\end{definition}

\begin{definition}[Contingent commodity]\label{def:b2:insurance-and-risk-sharing:contingent-commodity}
A \emph{contingent commodity}\index{contingent commodity} is a good delivered only if a given state
occurs: ``one dalen if it rains less than \qty{300}{mm} this season''. A plan of consumption under
risk is a bundle of contingent commodities, one per state.
\end{definition}

\begin{definition}[Arrow security]\label{def:b2:insurance-and-risk-sharing:arrow-security}
An \emph{Arrow security}\index{Arrow security} for state $s$ is an asset that pays one unit of money
if state $s$ occurs and nothing otherwise. With one Arrow security for every state, any pattern of
payments across states can be bought.
\end{definition}

The insurance contract of the previous section is an \omterm{def:b2:insurance-and-risk-sharing:arrow-security}{Arrow security} for the loss state, bought at the
price $\psi$ and paid for in every state. Plot wealth after no loss, $z_G$, across and wealth after a
loss, $z_B$, up. Buying cover moves the household from its endowment $(z_0, z_0 - \Lambda)$ along a
\emph{budget line} of slope $-(1 - \psi)/\psi$: giving up $\psi$ in the good state buys $1 - \psi$ net
in the bad one. The fair price gives the \emph{fair-odds line}, of slope $-(1 - \mu)/\mu$. The
\emph{certainty line} $z_G = z_B$ holds the riskless plans.

\begin{proposition}[Tangency on the certainty line]\label{prop:b2:insurance-and-risk-sharing:certainty-line}
Under expected utility, the marginal rate of substitution between wealth in the good and in the bad
state is $\mathrm{MRS} = (1 - \mu)u'(z_G)/(\mu\,u'(z_B))$. On the certainty line it equals $(1 -
\mu)/\mu$, the slope of the fair-odds line, for every utility of money.
\end{proposition}

\begin{proof}
Expected utility is $U = (1 - \mu)u(z_G) + \mu u(z_B)$, so along an indifference curve
\[
\frac{\dd z_B}{\dd z_G} = -\frac{(1 - \mu)\,u'(z_G)}{\mu\,u'(z_B)} .
\]
When $z_G = z_B$ the two marginal utilities cancel.
\end{proof}

So every risk-averse household's indifference curves cross the certainty line with the slope of the
fair odds, and a fair budget line is tangent there: this is
\cref{prop:b2:insurance-and-risk-sharing:mossin} again, in a picture. A loaded line is flatter, and the
tangency moves below the certainty line, towards the endowment.

\begin{omfigure}
\begin{tikzpicture}[econ]
\begin{axis}[omrecord, width=0.62\linewidth, height=0.58\linewidth, xmin=0, xmax=125, ymin=0,
    ymax=125, xlabel={Wealth with no loss, $z_G$ (thousands)}, ylabel={Wealth after a loss, $z_B$
    (thousands)}, grid=none, axis lines=left, xtick={40,80,100}, ytick={40,63.33,80},
    yticklabels={40,63,80}]
  \addplot[guide, domain=0:125, samples=2] {x};
  \addplot[omThm, thick, domain=60:120, samples=2] {40 - 2*(x - 100)};
  \addplot[omExo, thick, dashed, domain=60:120, samples=2] {40 - 1.5*(x - 100)};
  \addplot[curve, domain=58:125, samples=80] {80^3/x^2};
  \addplot[curve, thin, domain=58:125, samples=80] {63.333*84.444^2/x^2};
  \node[eqpoint, label={[font=\footnotesize]below left:$E$}] at (axis cs:100,40) {};
  \node[eqpoint, label={[font=\footnotesize]above left:$F$}] at (axis cs:80,80) {};
  \node[eqpoint, label={[font=\footnotesize]below left:$L$}] at (axis cs:84.444,63.333) {};
  \node[font=\footnotesize, anchor=west] at (axis cs:96,118) {certainty line};
  \node[font=\footnotesize, anchor=east, omThm] at (axis cs:59,120) {fair odds};
  \node[font=\footnotesize, anchor=east, omExo] at (axis cs:59,100) {loaded};
\end{axis}
\end{tikzpicture}

{\small Insurance in contingent goods: $z_0 = 100$, $\Lambda = 60$, $\mu = 1/3$ (thousands of dalens),
$u = \ln$. From the endowment $E$, the fair-odds line (slope $-2$) is tangent to an indifference curve
at $F$, full insurance at $80$. With a 20 per cent \omterm{def:b2:insurance-and-risk-sharing:loading}{loading} ($\psi = 0.4$, slope $-1.5$) the best point
is $L$: cover \num{38.9}, wealth $84.4$ and $63.3$.}
\end{omfigure}

\section{Efficient risk sharing}\label{sec:b2:insurance-and-risk-sharing:sharing}

Now take two households and no insurer. In state $s$ (probability $\mu_s$) household $i$ has the
endowment $\omega_i(s)$, and the aggregate endowment is $\Omega(s) = \omega_1(s) + \omega_2(s)$. An
allocation gives household $i$ the amounts $z_i(s)$ with $z_1(s) + z_2(s) = \Omega(s)$ in every
state.

\begin{definition}[Idiosyncratic and aggregate risk]\label{def:b2:insurance-and-risk-sharing:idiosyncratic-risk}
A household's \emph{idiosyncratic risk}\index{idiosyncratic risk} is the part of the variation of its
endowment that cancels out across households, leaving the aggregate unchanged.
\end{definition}

\begin{definition}[Aggregate risk]\label{def:b2:insurance-and-risk-sharing:aggregate-risk}
\emph{Aggregate risk}\index{aggregate risk} is variation of the aggregate endowment $\Omega(s)$ across
states. No arrangement among the households themselves can remove it: whatever they share, the total
they share is $\Omega(s)$.
\end{definition}

\begin{definition}[Efficient risk sharing]\label{def:b2:insurance-and-risk-sharing:efficient-risk-sharing}
An allocation achieves \emph{efficient risk sharing}\index{efficient risk sharing} if no other
allocation of the same aggregate endowment gives one household a higher expected utility without
giving the other a lower one.
\end{definition}

(\Cref{ch:b2:exchange-economies} gives this criterion its general form, for any goods and any number
of people.)

\begin{proposition}[The sharing rule (Borch)]\label{prop:b2:insurance-and-risk-sharing:borch}
Let $u_1$, $u_2$ be differentiable, increasing and strictly concave. An allocation with $z_i(s) > 0$
in every state is efficient only if the ratio of marginal utilities is the same in every state:
\[
\frac{u_1'(z_1(s))}{u_2'(z_2(s))} = \kappa \quad\text{for every } s,
\]
for some constant $\kappa > 0$.
\end{proposition}

\begin{proof}
Suppose the ratio is higher in state $s$ than in state $t$: $u_1'(s)/u_2'(s) > u_1'(t)/u_2'(t)$,
writing $u_i'(s)$ for $u_i'(z_i(s))$. Move a small amount $\epsilon > 0$ from household 2 to
household 1 in state $s$, and $\delta > 0$ from household 1 to household 2 in state $t$. To first
order the expected utilities change by
\[
\Delta U_1 = \mu_s u_1'(s)\,\epsilon - \mu_t u_1'(t)\,\delta, \qquad \Delta U_2 = -\mu_s
u_2'(s)\,\epsilon + \mu_t u_2'(t)\,\delta .
\]
Both are positive when $\frac{\mu_s u_2'(s)}{\mu_t u_2'(t)}\,\epsilon < \delta < \frac{\mu_s
u_1'(s)}{\mu_t u_1'(t)}\,\epsilon$, an interval that is not empty by the assumption. For $\epsilon$
small enough the first-order terms dominate, both households gain, and the allocation was not
efficient.
\end{proof}

\begin{corollary}[What efficient sharing does]\label{cor:b2:insurance-and-risk-sharing:aggregate-only}
In an efficient allocation with positive amounts:
\begin{enumerate}
  \item if there is no \omterm{def:b2:insurance-and-risk-sharing:aggregate-risk}{aggregate risk} ($\Omega(s)$ the same in every state), each household's
        consumption is the same in every state: full mutual insurance;
  \item in general, each household's consumption depends on the state only through $\Omega(s)$, and
        rises with it.
\end{enumerate}
\end{corollary}

\begin{proof}
Take two states $s$, $t$ with $\Omega(s) \ge \Omega(t)$, and suppose $z_1(s) < z_1(t)$. Then
$u_1'(s) > u_1'(t)$, so by \cref{prop:b2:insurance-and-risk-sharing:borch} $u_2'(s) > u_2'(t)$, that
is $z_2(s) < z_2(t)$. Adding up, $\Omega(s) < \Omega(t)$: a contradiction. So $\Omega(s) \ge
\Omega(t)$ implies $z_1(s) \ge z_1(t)$, and likewise $z_2(s) \ge z_2(t)$. With $\Omega(s) =
\Omega(t)$, applying this both ways gives equality, which proves (1) and the first half of (2).
\end{proof}

Efficient sharing ignores whose boat came back full: only the size of the total catch matters. Each
household's own endowment, once the aggregate is known, tells nothing about its consumption. This is
the prediction the tests of \cref{sec:b2:insurance-and-risk-sharing:tests} look for.

\begin{example}[Two CARA households]\label{ex:b2:insurance-and-risk-sharing:cara}
With $u_i(z) = -e^{-a_i z}/a_i$, $u_i'(z) = e^{-a_i z}$ and the rule says $a_2 z_2(s) - a_1 z_1(s)$ is
the same in every state. With $z_2 = \Omega - z_1$,
\[
z_1(s) = \frac{a_2}{a_1 + a_2}\,\Omega(s) + k,
\]
for a constant $k$ that sets how the expected total is divided. Each household bears a fixed share of
the \omterm{def:b2:insurance-and-risk-sharing:aggregate-risk}{aggregate risk}, its risk tolerance $1/a_i$ divided by the total risk tolerance: the less
risk-averse household bears more.
\end{example}

\begin{omfigure}
\begin{tikzpicture}[econ, scale=0.055]
  \draw[thick] (0,0) rectangle (100,100);
  \draw[guide] (0,0) -- (100,100);
  \node[below left] at (0,0) {$O_1$};
  \node[above right] at (100,100) {$O_2$};
  \node[below, font=\footnotesize] at (50,-1) {household 1's wealth in state $G$ $\rightarrow$};
  \node[above, font=\footnotesize] at (50,101) {$\leftarrow$ household 2's wealth in state $G$};
  \node[rotate=90, above, font=\footnotesize] at (-1,50) {household 1, state $B$ $\rightarrow$};
  \draw[omDef, thick, domain=16:100, samples=80] plot (\x, {1600/\x});
  \draw[omThm, thick, domain=16:100, samples=80] plot ({100 - \x}, {100 - 1600/\x});
  \draw[omDef, thin, domain=40:100, samples=60] plot (\x, {2500/\x});
  \draw[omThm, thin, domain=40:100, samples=60] plot ({100 - \x}, {100 - 2500/\x});
  \node[eqpoint, label={[font=\footnotesize]right:$\omega$}] at (80,20) {};
  \node[eqpoint] at (40,40) {};
  \node[eqpoint] at (60,60) {};
  \node[font=\footnotesize, anchor=west] at (62,74) {efficient allocations};
  \draw[->, thin] (66,72) -- (55,56);
\end{tikzpicture}

{\small Two households, two equally likely states $G$ and $B$, no \omterm{def:b2:insurance-and-risk-sharing:aggregate-risk}{aggregate risk}: the total is 100 in
both states, so the box is square. Household 1 is rich in $G$, household 2 in $B$ (endowment
$\omega$). Indifference curves (blue for household 1, red for household 2, here $u = \ln$) are tangent
only on the diagonal, which is both households' certainty line: every efficient allocation is full
mutual insurance.}
\end{omfigure}

Pooling is the same logic with many households.

\begin{proposition}[Pooling]\label{prop:b2:insurance-and-risk-sharing:pooling}
Let $n$ households each face an independent loss with mean $\mu\Lambda$ and variance $\sigma^2$, and
let them share the total equally. Each then bears a loss with the same mean $\mu\Lambda$ and variance
$\sigma^2/n$; by the small-risk approximation its \omterm{def:b2:expected-utility:risk-premium}{risk premium} falls from about
$\tfrac12\mathcal{A}\sigma^2$ to about $\tfrac12\mathcal{A}\sigma^2/n$.
\end{proposition}

\begin{proof}
Each bears $\frac1n\sum_j \tilde x_j$; the mean is $\mu\Lambda$ and, the losses being independent, the
variance is $n\sigma^2/n^2$.
\end{proof}

This is the insurer's business: by the law of large numbers, the average claim of a large pool is
almost certain, so the insurer can sell cover near the expected payment. Pooling fails when the losses
are not independent (a drought, a flood, a recession hit everyone at once), and then the premium of
the shared risk stays large however many join.

\section{A first portfolio choice}\label{sec:b2:insurance-and-risk-sharing:portfolio}

Buying insurance reduces risk at a price; buying a risky asset takes on risk for a reward. The same
tools answer both.

\begin{definition}[Portfolio share]\label{def:b2:insurance-and-risk-sharing:portfolio-share}
A household with wealth $z_0$ that puts the amount $\theta z_0$ in a risky asset and the rest in a safe
one holds the \emph{portfolio share}\index{portfolio share} $\theta$. Its wealth next period is $z_0(1
+ r + \theta\tilde\xi)$, where $r$ is the safe real rate and $\tilde\xi$ the risky asset's excess
return.
\end{definition}

\begin{proposition}[A small risky holding]\label{prop:b2:insurance-and-risk-sharing:arrow-portfolio}
A risk-averse expected-utility maximiser holds some of the risky asset ($\theta^{*} > 0$) if and only
if its expected excess return $\bar\xi$ is positive.
\end{proposition}

\begin{proof}
Let $V(\theta) = \mathbb{E}[u(z_0(1 + r + \theta\tilde\xi))]$, which is concave in $\theta$. Then
$V'(0) = z_0\,u'(z_0(1 + r))\,\bar\xi$, positive exactly when $\bar\xi > 0$. If $V'(0) > 0$ a small
holding raises $V$; if $V'(0) \le 0$, concavity makes $\theta = 0$ at least as good as every positive
share.
\end{proof}

So the risk-averse household that insures its house also buys some shares: at the margin, a small
risk costs a premium of second order (\cref{prop:b2:expected-utility:small-risk}) and earns a return of
first order.

\begin{proposition}[The share, small risks]\label{prop:b2:insurance-and-risk-sharing:share}
For small risks, $\theta^{*} \approx \bar\xi/(\mathcal{R}\,\sigma^2)$, with $\mathcal{R}$ the relative
\omterm{def:b2:expected-utility:risk-aversion}{risk aversion} at $z_0(1 + r)$. Under CRRA the share does not depend on wealth.
\end{proposition}

\begin{proof}
To second order, $u(z_0(1 + r) + z_0\theta\tilde\xi) \approx u + u'z_0\theta\tilde\xi +
\tfrac12u''z_0^2\theta^2\tilde\xi^2$. Taking expectations, and neglecting $\bar\xi^2$ against
$\sigma^2$, $V(\theta) \approx u + u'z_0\bigl(\theta\bar\xi - \tfrac12\mathcal{R}\theta^2\sigma^2
\bigr)$ with $\mathcal{R} = -z_0(1 + r)u''/u'$ and $1 + r \approx 1$. The bracket is largest at $\theta
= \bar\xi/(\mathcal{R}\sigma^2)$. Under CRRA, $\mathcal{R} = \gamma$ at every wealth.
\end{proof}

\begin{proposition}[CARA and normal returns]\label{prop:b2:insurance-and-risk-sharing:cara-normal}
If utility is CARA with coefficient $a$ and $\tilde\xi$ is normal, the amount invested in the risky
asset is exactly $\theta^{*}z_0 = \bar\xi/(a\sigma^2)$, whatever the wealth.
\end{proposition}

\begin{proof}
With $A = \theta z_0$ invested, final wealth is normal with mean $z_0(1 + r) + A\bar\xi$ and variance
$A^2\sigma^2$. For a normal $\tilde y$, $\mathbb{E}[e^{-a\tilde y}] = e^{-a\mathbb{E}[\tilde y] +
a^2\mathrm{Var}(\tilde y)/2}$, so maximising expected utility is maximising $A\bar\xi -
\tfrac12aA^2\sigma^2$, whose maximum is at $A = \bar\xi/(a\sigma^2)$.
\end{proof}

\begin{omfigure}
\begin{tikzpicture}[econ]
\begin{axis}[omrecord, width=0.66\linewidth, height=0.42\linewidth, xmin=1, xmax=10, ymin=0,
    ymax=1.6, xlabel={Relative risk aversion $\gamma$}, ylabel={Portfolio share $\theta^{*}$}, grid=none,
    xtick={1,2,3,4,5,6,7,8,9,10}]
  \addplot[omDef, thick, domain=1:10, samples=100] {0.05/(x*0.0324)};
  \addplot[guide, domain=1:10, samples=2] {1};
  \node[font=\footnotesize, anchor=south west] at (axis cs:5.5,1) {all wealth in the risky asset};
  \node[eqpoint] at (axis cs:3,0.5144) {};
  \node[font=\footnotesize, anchor=south west] at (axis cs:3.1,0.53) {$\gamma = 3$: $\theta^{*} = 0.51$};
\end{axis}
\end{tikzpicture}

{\small The \omterm{def:b2:insurance-and-risk-sharing:portfolio-share}{portfolio share} $\bar\xi/(\gamma\sigma^2)$ for an expected excess return of 5 per cent and a
standard deviation of 18 per cent, round numbers chosen for the illustration. Above the dotted line the
household borrows to invest. \Cref{ch:b2:asset-prices} confronts such numbers with the returns
actually earned.}
\end{omfigure}

\section{Testing full risk sharing}\label{sec:b2:insurance-and-risk-sharing:tests}

With many households and identical attitudes to risk, efficient sharing has a sharp implication.

\begin{proposition}[Constant shares]\label{prop:b2:insurance-and-risk-sharing:constant-shares}
Let $n$ households have the same CRRA utility with coefficient $\gamma$. In an efficient allocation
with positive amounts, each household consumes a constant share of the aggregate: $c_i(s) =
k_i\,\Omega(s)$, so that $\ln c_i(s) - \ln \bar c(s)$ does not depend on $s$, where $\bar c(s)$ is the
geometric mean of the households' consumptions.
\end{proposition}

\begin{proof}
The argument of \cref{prop:b2:insurance-and-risk-sharing:borch} applies to any pair $i$, $j$:
$u'(c_i(s))/u'(c_j(s)) = (c_i(s)/c_j(s))^{-\gamma}$ is the same in every state, so $c_i(s)/c_j(s)$ is.
Then $c_i(s)/\Omega(s) = c_i(s)/\sum_j c_j(s) = 1/\sum_j (c_j(s)/c_i(s))$ is constant. The log of
$c_i/\bar c$ is the log of a ratio of constants.
\end{proof}

Year to year, then, every household's consumption grows at the same rate as the group's, whatever
happens to its own income. The test regresses the growth of household consumption on the growth of
average consumption and of the household's own income.

\begin{measure}[The consumption-insurance regression]
With a panel of households $i$ observed over years $t$, estimate
\[
\Delta\ln c_{it} = \beta_0 + \beta_1\,\Delta\ln\bar c_t + \beta_2\,\Delta\ln m_{it} + u_{it},
\]
where $\bar c_t$ is the group's average consumption (the village, the region, the country) and
$m_{it}$ the household's income. Full sharing predicts $\beta_1 = 1$ and $\beta_2 = 0$: own income
matters only through the aggregate, which $\bar c_t$ already measures. Under partial sharing, in which
a household keeps a share of its own shock, $\beta_2$ estimates that share. Two cautions.
Measurement error in consumption only adds noise; measurement error in income biases $\hat\beta_2$
towards zero and makes sharing look better than it is. And income can respond to choices (work
more after a bad harvest), so the test is strongest with shocks the household does not choose:
illness, rainfall on its own plots, a lay-off.
\end{measure}

\begin{model}{Full risk sharing}
\assumes Households in a group can make and enforce any agreement to share their incomes across
states; they maximise expected utility with the same CRRA coefficient.
\predicts Household consumption moves one for one with the group's and does not respond to own
income once the group's consumption is controlled for ($\beta_1 = 1$, $\beta_2 = 0$).
\tested Village panels in India (Townsend 1994); national accounts across countries (this chapter's
two charts).
\breaks Everywhere at the margin: own income keeps a small but significant coefficient in villages,
and across countries consumption follows domestic output far more than world consumption.
\end{model}

\begin{record}[Village risk sharing]
Townsend (1994) followed about 40 households in each of three villages of semi-arid southern India,
Aurepalle, Shirapur and Kanzara, from 1975 to 1984.
% ledger: b2:townsend1994-design
Harvests there are risky, and formal insurance was absent. Household by household, he regressed
consumption on village average consumption and on own income. The average coefficient on village
consumption was $0.97$, $0.94$ and $0.47$ in the three villages; the average coefficient on own income
% ledger: b2:townsend1994-cbar
was $0.11$, $-0.05$ and $0.11$, and the averages for every measure of income stayed below $0.35$ but
% ledger: b2:townsend1994-income
one. Full sharing is rejected statistically, yet, in the author's words, it gives a surprisingly good
benchmark: in villages with no insurance company, families, neighbours, lenders and grain stores
spread most of the \omterm{def:b2:insurance-and-risk-sharing:idiosyncratic-risk}{idiosyncratic risk}. The landless were less well insured than their neighbours in
one of the three villages.
\end{record}

Across countries the same test fails much more clearly. With full international sharing and identical
CRRA utilities, every country's consumption would grow at the world's rate. Its consumption growth
would correlate perfectly with world consumption growth, and its volatility would be the world's,
1.79 points, below the GDP volatility of all but 8 of the 123 countries. The opening chart shows the
% figdata: summary (world_sd_cons 1.79, n_sd_gdp_below_world 8)
opposite. The chart below checks the correlations: in 71 per cent of countries, consumption growth is
% figdata: summary (share_cons_corr_below_gdp 70.7, median_corr_gdp 0.54, median_corr_cons 0.38)
less correlated with the world's consumption growth than output growth is with the world's output
growth (medians $0.38$ and $0.54$).

\begin{omchart}{Correlation of each country's growth with world growth, 1991--2023: GDP per person
  (across) and household consumption per person (up). Full international risk sharing puts every point
  at the top of the chart; points below the dotted line are countries whose consumption follows the
  world less than their output does.}{World Bank, World Development Indicators: NE.CON.PRVT.PC.KD.ZG
  and NY.GDP.PCAP.KD.ZG, each country against the World aggregate of the same indicator (which
  includes the country itself)}
\begin{tikzpicture}
\begin{axis}[omscatter, xmin=-0.4, xmax=1, ymin=-0.4, ymax=1, width=0.68\linewidth,
    height=0.6\linewidth, xlabel={Correlation of GDP growth with the world's},
    ylabel={Correlation of consumption growth with the world's}]
  \addplot+[mark size=1.3pt] table[x=corr_gdp, y=corr_cons]
    {figdata/bachelor-2/02-insurance-and-risk-sharing/out/correlation.csv};
  \draw[black!70, densely dotted, thick] (axis cs:-0.4,-0.4) -- (axis cs:1,1);
\end{axis}
\end{tikzpicture}
\end{omchart}

Part of the gap may be measurement. Private consumption leaves out public services; much of what
households consume (housing, haircuts) is not traded across borders, so its price moves with local
conditions; and consumption, like output, is measured with error. The model points to two other
reasons, each a claim to test rather than a finding of this chart: a household that holds most of its
wealth at home shares little of its country's risk, and contracts across borders are harder to enforce
than village ones, where the neighbour who does not pay back can be excluded from the next round of
help. \Cref{exo:b2:insurance-and-risk-sharing:9} asks which observation would tell these apart.

\section{Answers outside insurance markets}\label{sec:b2:insurance-and-risk-sharing:outside}

Formal insurance is one answer among several, and each has its evidence and its limits.

\begin{itemize}
  \item \textbf{Families and neighbours.} The village record shows informal sharing doing most of the
        work where no insurer operates. Its limit is the one \cref{prop:b2:insurance-and-risk-sharing:pooling}
        names: a village shares its \omterm{def:b2:insurance-and-risk-sharing:idiosyncratic-risk}{idiosyncratic risk}, but a drought hits all its fields together.
  \item \textbf{Mutual societies and cooperatives.} Members insure each other and own the pool, so no
        profit margin is charged; they still face the information problems of
        \cref{ch:b2:adverse-selection,ch:b2:moral-hazard}, and a small pool carries its members' common
        risks. The weekend problem prices one.
  \item \textbf{Index insurance.} A contract that pays on a measured index (rainfall at the nearest
        station) rather than on the farmer's own loss costs little to settle and cannot be
        manipulated by the farmer, but it pays when the index says so, not when the field fails: the
        gap is called basis risk.
  \item \textbf{Public insurance.} A state programme can pool the whole population and make
        membership compulsory, which solves the selection problem of \cref{ch:b2:adverse-selection}.
        Its prices are set by law rather than by claims, so some members pay more than their expected
        cost and others less, and subsidised cover changes where people build and what they plant;
        government failure is a cost to weigh, as market failure is.
\end{itemize}

\begin{record}[Index insurance for farmers]
Karlan et al.\ (2014) offered rainfall index insurance to farmers in northern Ghana at prices drawn at
random. In one year, of \num{1095} farmers offered cover, 63.9 per cent bought at a subsidised price,
% ledger: b2:karlan2014-demand
55.6 per cent at the actuarially fair price and 40.0 per cent at the market price. Across their
experiments, 40 to 50 per cent bought at the fair price, and buyers covered more than 60 per cent of
% ledger: b2:karlan2014-acreage
their acreage. Insured farmers spent \usd{266} more on cultivation (standard error \usd{134}) than the
% ledger: b2:karlan2014-invest
control group, whose mean was \usd{2058}: uninsured risk, more than lack of credit, held back their
investment. Demand the next year rose after the farmer, or others in his network, had received a
payout, consistent with basis risk and with imperfect trust that payouts would be made. The model of
\cref{sec:b2:insurance-and-risk-sharing:demand} gets the direction of the price effect right; it does
not predict that only about half the farmers buy at a fair price.
\end{record}

\section{Exercises}\label{sec:b2:insurance-and-risk-sharing:exercises}

\begin{exercise}[$\star$]\label{exo:b2:insurance-and-risk-sharing:1}
In the table behind the opening chart, the standard deviations of GDP and consumption growth are 1.76
% figdata: volatility (USA 1.758 1.966; JPN 2.05 1.503)
and 1.97 points for the United States, 2.05 and 1.50 for Japan. Place each country with respect to
the dotted line. Under full international risk sharing with identical CRRA utilities, what would each
country's consumption volatility be, and where would the points lie?
\end{exercise}

\begin{exercise}[$\star$]\label{exo:b2:insurance-and-risk-sharing:2}
A house worth \dal{40000} burns down in a year with probability $0.005$. Find the \omterm{def:b2:insurance-and-risk-sharing:fair-premium}{actuarially fair
premium} for full cover. An insurer charges \dal{260}: find the \omterm{def:b2:insurance-and-risk-sharing:loading}{loading}.
\end{exercise}

\begin{exercise}[$\star$]\label{exo:b2:insurance-and-risk-sharing:3}
A household has \dal{50000} and loses \dal{20000} with probability $0.1$; cover costs $\psi = 0.12$ per
dalen. Write $z_G$ and $z_B$ as functions of the cover. Find the slope of the budget line, the slope of
the fair-odds line, the \omterm{def:b2:insurance-and-risk-sharing:loading}{loading}, and the point where the budget line meets the certainty line.
\end{exercise}

\begin{exercise}[$\star$]\label{exo:b2:insurance-and-risk-sharing:4}
A survey follows households of a village for ten years. Name two observations that would reject full
risk sharing within the village, and one that would not, even though it may look like it (think of a
drought).
\end{exercise}

\begin{exercise}[$\star\star$]\label{exo:b2:insurance-and-risk-sharing:5}
For the household of \cref{exo:b2:insurance-and-risk-sharing:3}, with $u = \ln$, find the best cover
and its wealth in each state. Check that it is below full cover, as
\cref{prop:b2:insurance-and-risk-sharing:mossin} says.
\end{exercise}

\begin{exercise}[$\star\star$]\label{exo:b2:insurance-and-risk-sharing:6}
Predict, then compute: what happens to the household's best cover if the price rises to $\psi = 0.15$?
Find the price, and the \omterm{def:b2:insurance-and-risk-sharing:loading}{loading}, at which it stops buying cover altogether.
\end{exercise}

\begin{exercise}[$\star\star$]\label{exo:b2:insurance-and-risk-sharing:7}
Two Brenland farmers have CARA utilities with $a_1 = \num{0.0002}$ and $a_2 = \num{0.0001}$ per dalen.
Their joint harvest is \dal{100000} or \dal{60000}, with probability $\tfrac12$ each. Find the
efficient sharing rule that gives farmer 1 an expected income of \dal{30000}, and each farmer's
income in each state. Which farmer bears more of the \omterm{def:b2:insurance-and-risk-sharing:aggregate-risk}{aggregate risk}, and why?
\end{exercise}

\begin{exercise}[$\star\star$]\label{exo:b2:insurance-and-risk-sharing:8}
With the World Bank series of the opening chart (a few lines of Python with pandas do it), compute, for
1991--2023, the standard deviations of the two growth rates and their ratio for the United States,
Japan, Korea, France, Brazil and India. Which of these countries come closest to the prediction of full
international sharing?
\end{exercise}

\begin{exercise}[$\star\star$]\label{exo:b2:insurance-and-risk-sharing:9}
Read the chart of correlations. What does full international sharing predict for every point? Give
two reasons why measured correlations could fall short of that prediction even if financial markets
shared risk perfectly, and one observation that would separate those reasons from a genuine lack of
sharing.
\end{exercise}

\begin{exercise}[$\star\star\star$]\label{exo:b2:insurance-and-risk-sharing:10}
Two households have CRRA utilities with different coefficients $\gamma_1 < \gamma_2$. Show that in an
efficient allocation $\gamma_1\ln c_1(s) - \gamma_2\ln c_2(s)$ is the same in every state, and deduce
that household 1's share of the aggregate rises with the aggregate. What would the regression of the
measure box find for household 1 when the aggregate grows, and is full sharing then rejected by a
coefficient $\beta_1 \neq 1$?
\end{exercise}

\begin{exercise}[$\star\star\star$]\label{exo:b2:insurance-and-risk-sharing:11}
$n$ households with CARA utility, $a = \num{0.0001}$ per dalen, each lose \dal{10000} with probability
$0.1$, independently, and share the total loss equally. Show that the certainty-equivalent cost of the
risk to each household is $\frac{n}{a}\ln\bigl(1 - \mu + \mu e^{a\Lambda/n}\bigr)$. Evaluate it for $n
= 1$, $10$ and $100$, prove that it falls with $n$ towards the expected loss, and say what happens if
the losses are perfectly correlated.
\end{exercise}

\begin{exercise}[$\star\star\star$]\label{exo:b2:insurance-and-risk-sharing:12}
With an expected excess return of 5 per cent and a standard deviation of 18 per cent, find the
\omterm{def:b2:insurance-and-risk-sharing:portfolio-share}{portfolio share} of a CRRA household with $\gamma = 3$. Under CRRA, and under CARA with normal returns,
what happens to the share and to the amount invested when the household's wealth doubles? Which
pattern, in data on the portfolios of rich and poor households, would reject each model?
\end{exercise}

\section{Problem: The Kestrian fishing cooperative}\label{sec:b2:insurance-and-risk-sharing:problem}

\begin{problem}[{Weekend problem --- two boats of Kestria, an invented economy, whose catches are
risky and correlated: sharing within a cooperative, insurance bought outside, and a survey that tests
how much risk the cooperative really shares}]
\label{pb:b2:insurance-and-risk-sharing:1}
Two fishing households each own a boat. A boat's catch earns \kst{30000} in a good season and
\kst{10000} in a bad one. The four states have probabilities: both good $0.4$, only boat 1 good $0.2$,
only boat 2 good $0.2$, both bad $0.2$. Each household maximises $\mathbb{E}[\ln z]$ of its income.
Round amounts to the nearest kest.

\medskip
\noindent\textbf{Part I --- Two risky catches.}
\begin{enumerate}
  \item Find the probability that a boat has a good season and that both have one. Are the catches
        independent?
  \item Find the mean and the standard deviation of one household's income.
  \item Find the correlation coefficient between the two incomes.
  \item Give the aggregate catch in each state and its distribution.
  \item Find the expected utility and the \omterm{def:b2:expected-utility:certainty-equivalent}{certainty equivalent} of a household that keeps its own catch.
  \item Write each household's income as half the aggregate plus a remainder. Show that the two
        remainders add up to zero in every state, and find the standard deviation of each part.
\end{enumerate}

\medskip
\noindent\textbf{Part II --- The cooperative.}
\begin{enumerate}[resume]
  \item Use \cref{prop:b2:insurance-and-risk-sharing:constant-shares} to show that every efficient
        allocation gives household 1 a constant share $k$ of the aggregate catch.
  \item The cooperative splits the catch equally. Give each household's income in each state, its mean,
        its expected utility and its \omterm{def:b2:expected-utility:certainty-equivalent}{certainty equivalent}.
  \item What is the cooperative worth to each household, in kests a year?
  \item By what percentage does the variance of a household's income fall? What share of its
        \omterm{def:b2:insurance-and-risk-sharing:idiosyncratic-risk}{idiosyncratic risk} is removed?
  \item Suppose instead that both skippers have CARA utility and that skipper 1 is twice as risk-averse
        as skipper 2. Find the efficient rule that keeps skipper 1's expected income at \kst{22000},
        and his income in each state.
  \item Why can the cooperative not remove the risk that remains, and who could take it?
\end{enumerate}

\medskip
\noindent\textbf{Part III --- An insurer outside.}
An insurer sells cover against a household's own bad season, at a \omterm{def:b2:insurance-and-risk-sharing:loading}{loading} of 20 per cent.
\begin{enumerate}[resume]
  \item Find the fair price and the loaded price per kest of cover.
  \item A household on its own (no cooperative) buys cover against a fall of \kst{20000}. Find its best
        cover, using the formula of \cref{ex:b2:insurance-and-risk-sharing:loading}.
  \item Find its income in each state and its \omterm{def:b2:expected-utility:certainty-equivalent}{certainty equivalent}.
  \item Compare with the cooperative. Which arrangement does the household prefer, and why, although
        the insurer charges a \omterm{def:b2:insurance-and-risk-sharing:loading}{loading} and the cooperative does not? What does the insurer expect to
        earn on the contract?
  \item The cooperative shares the catch equally and buys, for each household, cover paid only when
        both seasons are bad, at a \omterm{def:b2:insurance-and-risk-sharing:loading}{loading} of 20 per cent. Find the best cover (numerically) and the
        \omterm{def:b2:expected-utility:certainty-equivalent}{certainty equivalent}.
  \item Rank the four arrangements of Parts I--III. What would a cover that pays on a household's own
        bad season do to the care its skipper takes, and does sharing in the cooperative have the same
        effect?
\end{enumerate}

\medskip
\noindent\textbf{Part IV --- The cooperative's survey.}
A larger Kestrian cooperative of 100 households records their incomes and consumptions for 8 years.
The regression of the measure box on the 700 yearly changes gives:

\begin{center}
\pgfplotstabletypeset[col sep=comma,
  columns={term,coef,se,t},
  columns/term/.style={string type, column name={Regressor}, column type={l},
    string replace={const}{constant},
    string replace={dlncbar}{$\Delta\ln\bar c_t$},
    string replace={dlnm}{$\Delta\ln m_{it}$}},
  columns/coef/.style={column name={Coefficient}, column type={r},
    fixed, fixed zerofill, precision=3},
  columns/se/.style={column name={Standard error}, column type={r},
    fixed, fixed zerofill, precision=3},
  columns/t/.style={column name={$t$}, column type={r}, fixed, fixed zerofill, precision=1},
  every head row/.style={before row=\toprule, after row=\midrule},
  every last row/.style={after row=\bottomrule}]
  {figdata/bachelor-2/02-insurance-and-risk-sharing/out/coop_regression.csv}
\end{center}
{\small (The constant is $\beta_0$. The survey is simulated.)}
\begin{enumerate}[resume]
  \item What values of $\beta_1$ and $\beta_2$ does full sharing predict, and why?
  \item Read $\hat\beta_1$ and $\hat\beta_2$ with their standard errors.
  \item Test $\beta_2 = 0$ at the 5 per cent level.
  \item Test $\beta_1 = 1$. Compute $\hat\beta_1 + \hat\beta_2$, and explain why a household that keeps a
        share of its own shock produces estimates that add up to about one.
  \item What share of the idiosyncratic income risk does the cooperative insure?
  \item If incomes were recorded with error, in which direction would $\hat\beta_2$ move, and what
        would that do to the answer to question 23?
  \item State the result: what the two-boat cooperative is worth to a household in kests a year, and
        what share of \omterm{def:b2:insurance-and-risk-sharing:idiosyncratic-risk}{idiosyncratic risk} the larger cooperative insures.
\end{enumerate}
\end{problem}
